\section{Неполная наблюдаемость и информационные режимы}

\subsection{Латентное состояние и наблюдаемые сигналы}

Центральный банк не наблюдает \(q_t\) напрямую. Наблюдения задаются как шумные сигналы
\[
o_t=M_\varphi q_t+\nu_t,
\]
где \(\varphi\) обозначает информационный режим. Матрица \(M_\varphi\) определяет, какие
компоненты состояния доступны наблюдению, а \(\nu_t\) описывает ошибку измерения. Информационное
множество регулятора:
\[
\mathcal I_t^\varphi
=
\sigma(o_0^\varphi,\ldots,o_t^\varphi,i_0,\ldots,i_{t-1}).
\]
Информационное состояние
\[
s_t^\varphi=\Phi_\varphi(\mathcal I_t^\varphi)
\]
является сжатием истории наблюдений для выбора ставки. В дальнейшем именно \(s_t^\varphi\), а не
полная история \(\mathcal I_t^\varphi\), подаётся на вход правилу процентной ставки. Поэтому
разные режимы \(\varphi\) задают альтернативные способы наблюдать и сжимать одну и ту же
локальную экономическую динамику.

\subsection{Агрегатная оценка состояния}

Основная база сравнения в работе -- агрегатная оценка состояния по фильтру Калмана:
\[
s_t^{agg}
=
\mathbb E
\left[
(\pi_t,y_t,C_t)\mid \mathcal I_t^{agg}
\right].
\]
Она служит базой, относительно которой измеряется дополнительная ценность распределительных
сигналов. Если положительный эффект распределительных сигналов возникает только из-за ошибок
измерения агрегатов, он должен исчезать или заметно ослабевать после перехода к агрегатной оценке
состояния.

\subsection{Сопоставление информационных режимов}

В эксперименте сравниваются следующие информационные режимы:
\[
s_t^{obs\text{-}agg}=(\pi_t^{obs},y_t^{obs},C_t^{obs}),
\]
\[
s_t^{hist\text{-}agg}=(o_t^{agg},o_{t-1}^{agg},\ldots,i_{t-1}),
\]
\[
s_t^{obs\text{-}dist}
=
(\pi_t^{obs},y_t^{obs},C_t^{obs},
\overline{\kappa}_t^{obs},\ell_t^{obs},\xi_t^{r,obs}),
\]
\[
s_t^{agg+dist}
=
\mathbb E
\left[
(\pi_t,y_t,C_t,\overline{\kappa}_t,\ell_t,\xi_t^r)
\mid
\mathcal I_t^{dist}
\right],
\]
а также варианты, где к агрегатной оценке состояния добавляется только одна распределительная
статистика. Состояние \(s_t^{obs\text{-}dist}\) используется как дополнительный режим
сопоставления: оно показывает значение текущих агрегатных наблюдений вместе с наблюдаемыми
распределительными сигналами без полного восстановления распределительного состояния. Этот режим
не является основной базой сравнения, поскольку одновременно меняет состав наблюдений и способ
построения информационного состояния.

Режим наблюдаемого локального состояния
\[
s_t^{full}=q_t
\]
служит эталоном полной информации и задаёт верхнюю границу качества в рассматриваемой линейной
задаче.

Распределительные статистики и распределительные сигналы различаются. Статистики --
\(\overline{\kappa}_t\), \(\ell_t\), \(\xi_t^r\) -- являются экономическими объектами, построенными
из распределения домохозяйств. Сигналы -- это шумные или фильтрованные наблюдения этих объектов.
Слово ``признаки'' используется ниже только в регрессионных и прогностических проверках.
